\section{Introduction}\label{sec:introduction}

Smart contracts are software components that
run inside a blockchain. Their execution
modifies the state of the application implemented by the blockchain.
For instance, Ethereum smart contracts are typically written in
Solidity, can be installed in the state of the application and their
execution modifies a key/value store, used to represent accounts, balances
and actual data objects~\citep{AntonopoulosWPMP25}. Smart contracts
deal with money, notarization and agreements. They are immutable after
deployment. Therefore, their correctness is important and there exist
many techniques to increase confidence on their correctness.
The simpler and most traditional one is testing, which consists
in creating some executable usage scenarios for the smart contracts,
as a sequence of transactions starting from an initial fixture.
Testcases test the final state of the application to verify the correct
behavior of the smart contracts.

\begin{figure*}[t]
  \begin{center}
    \begin{tikzpicture}
      \draw [fill=verylightviolet,thick] (-0.5,-1.5) rectangle +(4,3);
      \draw (1.5,0.8) node {blockchain peer};
      \draw (1.5,0.4) node {{\small (networking and consensus)}};
      \draw[-, thick] (1.5,-1.5) -- (1.5,-2.7);
      \draw (1.5,-2.7) node[scale=0.7] {\db};
      \draw (1.5,-3.6) node {database of blocks};
      \draw[<->, thick] (3.5,-1) -- (5.5,-1);
      \draw (7.1,-1) node {other blockchain peers};
      \draw (4.5,-1.2) node {{\scriptsize gossiping}};

      \draw [fill=verylightyellow,thick] (0.5,-0.7) rectangle +(2,0.5);
      \draw (1.5,-0.45) node {mempool};
      \draw[->, thick] (-4,-0.45) -- (0.5,-0.45);
      \draw (-2.25,-0.7) node {transaction requests};
      \draw (-2.25,-0.2) node {{\small API}};
      \draw[<-, thick] (-4,0) -- (-0.5,0);
      \draw (-2.25,0.3) node {transaction responses};
      \draw [fill=verylightblue,thick] (-6,-0.6) rectangle +(2,0.8);
      \draw (-5,-0.2) node {client(s)};

      \draw [fill=verylightred,thick] (6,0.35) rectangle +(2,0.8);
      \draw (7.0,0.75) node {miner(s)};

      \draw [fill=verylightgreen,thick] (0.5,3) rectangle +(2,0.8);
      \draw (1.5,3.4) node {application};
      \draw[-, thick] (0.5,3.4) -- (-1,3.4);
      \draw (-1,3.4) node[scale=0.7] {\db};
      \draw (-1,2.5) node {database of states};

      \draw[->, thick] (3.5,1.0) -- (6,1.0);
      \draw (4.75,1.2) node {{\small challenges}};

      \draw[<-, thick] (3.5,0.5) -- (6,0.5);
      \draw (4.75,0.25) node {{\small answers}};

      \draw[<-, thick] (1.2,1.5) -- (1.2,3);
      \draw[<-, thick] (1.8,3) -- (1.8,1.5);
    \end{tikzpicture}
  \end{center}
  \caption{The architecture of a blockchain peer.
    Peer, miner(s) and application are separate software
    components, possibly running in distinct machines connected
    through a network. Client(s) are external
    software tools that connect to the peer through an API
    and run transactions to install smart contracts
    and run their methods.}\label{fig:architecture}
\end{figure*}

\Cref{fig:architecture} shows an abstract view of the architecture of
a typical blockchain peer. The peer deals with networking and consensus.
It communicates with one or more miners, that help the peer create new blocks,
stored in its database of blocks. Miners receive a challenge to solve
(for instance, the maximization of the hash of a new block) and provide
their best solution for the challenge. Peers are connected to other peers
and gossip new blocks and new transaction requests to them.
Clients are software
components, external to the peer, that consume the API of the peer.
For instance, they can be wallets that require the execution of money
transfer transactions. Or they can be DApps that require the execution
of transactions to install and interact with smart contracts in blockchain.
Finally, the application is a software component that specifies which
transactions are valid and which is their semantics. For instance, the Bitcoin
application has only transactions that transform input UTXOs into ouptut
UTXOs, while the Ethereum application has transactions that can execute
arbitrary EVM bytecode. Transaction requests are first parked inside a
mempool, a queue from where they are eventually extracted for execution.
Clients may receive responses about the outcome of the
transactions that they requested,
in the form of actual results or events. There is an actual API
that specifies the request/response
programming interface for the blockchain, such as the Web3 API,
in the case of Ethereum.

Testcases can be seen as clients that use the blockchain API to simulate
usage scenarios for smart contracts and verify their outcome.
In terms of \Cref{fig:architecture}, this means that testcases assume
the blockchain peer, its miners and its application to be working correctly
and use them as a testbench for the execution of the smart contracts, that
are the target of the testing activity. But another perspective is possible:
testcases can be assumed to pass and can get used as non-regression tests
for modifications to the blockchain peer, miners or application, or for
complete rewritings of the peer. The target of the testing activity are
peer, miners and application in this case.

This article uses testcases in poth ways.
It originated during the development of Hotmoka~\citep{hotmoka},
a new, permissionless, public blockchain based on Solidity-like
smart contracts written in a subset of Java called Takamaka~\citep{Spoto19}.
Testcases have been used to give confidence in the Takamaka
smart contracts, but also to support the development of the Hotmoka peer,
in order to guarantee that the latter executed the smart contracts
as expected. Moreover, the Hotmoka peer transitioned from an initial
version based on proof of stake over Tendermint~\citep{Kwon14} to its
current version based on proof of space over Mokamint~\citep{mokamint}.
Testcases on smart contracts have been invaluable during this transition,
to give confidence in the correctness of the new peer, despite its
radical shift of consensus algorithm.

This has been possible because the abstract API between clients and peer
is the same, regardless of the implementation of the peer and of its consensus
algorithm. The existence of this abstract interface has also allowed us
to apply the traditional stubbing approach to it: the peer
in \cref{fig:architecture} has been replaced with a stubbed peer that does
not implement consensus nor gossiping, but just forwards the transaction
requests to an actual application. The advantage is that
testcases can be executed much faster (minutes) than against an actual peer
of an actual blockchain (hours),
that must abide to consensus and creates blocks slowly, at regular intervals.
The limit is that stubbing only allows one to test the smart contract,
the application
(\cref{fig:architecture}) but not the full peer, since some components are mocked.
Nevertheless, this approach has been very useful to reduce the time
for the development of the Hotmoka application and currently runs during its
continuous integration as a Github action.

The contributions of this article are the following:
%
\begin{itemize}
\item the description of the abstract API between the clients and an
  Hotmoka blockchain peer, meant to allow the execution of object-oriented
  code in blockchain;
\item the description of how that API allows one to
  run testcases of smart contracts
  against different implementations of the peer, with different
  consensus algorithms;
\item a stubbed implementation of the peers, that makes the execution of the
  testcases faster and allows their use for continuous integration, but still keeps
  their same abstract API;
\item an adaptor for remote peers, that allows one to execute testcases
  against a local peer as well as against a remote peer, always using the
  same abstract API.
\end{itemize}

The use of a single API, for local, remote, actual or stubbed peers, allows one
to run the same testcases against all possible implementations, without modification,
in accordance with the open-closed principle~\citep{Meyer97}.

The rest of this article is organized as follows. \Cref{sec:related} reports
related works. \Cref{sec:takamaka} presents an example of smart contract
in the Takamaka subset of Java used by Hotmoka and of a testcase for that
smart contract. \Cref{sec:interface} presents the API between clients
and Hotmoka peers. \Cref{sec:stubbing} describes the implementation
of a stubbed Hotmoka peer.
\Cref{sec:adapting} shows how the API between clients and peers can be adapted to
a remote peer, potentially running in another machine and reachable thorugh
network calls. \Cref{sec:experiments} reports experiments with
the testing of Takamaka smart contracts against different implementations
of the API, including the stubbed one, both local and remote.
\Cref{sec:conclusion} concludes.
